% !TEX root = ../main.tex

\chapter{Conclusions}

In this thesis, we have shown a new generalisation of Phoa principle that applies on the initial $L$-algebra. This generalisation suggests a possible connection between the Segal completeness and the chain completeness. We expect that this hidden connection will be made explicit in future and will eventually guide us on constructing the models for synthetic domain theory.

Nevertheless, we should admit that the route which brings us from the classical Phoa principle to the transfinite Phoa principle heavily relies on the lattice structure of the interval type. As a result, Conjecture \ref{hypo-last} may see a non-negligible gap on types other than the observational algebras. One may expect a duality between the ``ordinary'' types and their observational algebras analogous to the Stone duality between the sober spaces and their observational algebras of open sets. One may ask if applying the functor $\Op$ twice gives such duality. It turns out that, if we define the a sober space to be a space $X$ satisfying $X\cong\Op^2(X)$, none of the simplices (either finite or transfinite) are sober, though we have the intuition that simplices contain the most points among all spaces capturing the same order structure. We hereby list two questions that remain to answer:
\begin{enumerate}
    \item Does the observational algebra $\Op(X):=X\to\I$ really behave like a synthetic counterpart of the frame of open sets, even though it does not guarantee arbitrary unions?
    \item Is there a general way to assign each ``good'' type (Segal, Rezk, $\I$-separated, etc.) a lattice that captures the information order and in reverse, reconstruct a type from the observational algebra?
\end{enumerate}
Sterling and Ye proposed another duality between spaces and lattices \cite{sterling_domains_2025}. They see $\Op$ as a functor sending an ``affine space'' to $\I$-algebras, and see the right adjoint of this functor as finding the spectrum of a $\I$-algebra. This framework might lead to a possible way to prove or disprove Conjecture \ref{hypo-last}.

This thesis also provides formalisations of the proofs in Cubical Agda. The formalisation codebase includes the following results:
\begin{enumerate}
    \item Axiomatisation of the lattice and $L$-algebra on $\I$.
    \item Relation between $L$-algebras and simplices.
    \item Classical and higher Phoa principles and their variants on spines.
    \item Transfinite Phoa principle and its variant on $\omega$-spine.
\end{enumerate}
We expect that the explicit constructions in our formalisation could contribute as a foundation for the future development of a full-scale calculus for synthetic domain theory.

Kudasov et al. have delivered a proof-assistant call Rzk \cite{kudasov_formalizing_2023} based on the synthetic $\infty$-category theory by Riehl et al. \cite{riehl_type_2023}, and it is promising to migrate our work into Rzk to utilize its native support for directed homotopy. Another work that provides a formalisation framework for synthetic domain theory is the guarded mode in Agda, based on the theory in \cite{mogelberg_bisimulation_2019} and \cite{kristensen_greatest_2022}. Despite the fact that the guarded mode is designed for the synthetic guarded domain theory \cite{birkedal_first_2012} instead of SDT mentioned in this thesis, we still see the possibility of implementing SDT as a language extension of Agda in the future.

% One crucial gap towards a full-scale calculus of SDT is the encoding of Axiom \ref{ax-inductive}. Forcing the initial $L$-algebra to be inductive by postulating this axiom directly may lead to potential difficulties in normalisation of the calculus, and hence developing from the language may be a possible way to tackle this issue.
